% Appendix table — pre-crisis comparison period (2022Q2).
% Generated by scripts/rq1_placebo.py. Do not hand-edit.
\begin{table}[htbp]
\centering
\caption{The uninsured-deposit share and deposit growth in the pre-crisis comparison period (2022Q2) and the banking-stress period (2023Q1)}
\label{tab:placebo_precrisis}
\begin{tabular}{lccc}
\toprule
 & \multicolumn{2}{c}{Period-specific OLS} & Pooled \\
\cmidrule(lr){2-3}\cmidrule(lr){4-4}
 & Pre-crisis & Banking stress & Two-period \\
 & (2022Q2) & (2023Q1) & model \\
\midrule
Uninsured share & -0.00898 & -0.07713 & -0.00592 \\
 & (0.01357) & (0.01927) & \\
 & [0.5084] & [0.0001] & [0.7780] \\
\addlinespace
Uninsured share $\times\ \mathbb{1}$[2023Q1] & & & \textbf{-0.06671} \\
 & & & (0.03995) \\
 & & & [\textbf{0.0953}] \\
\addlinespace
$\mathbb{1}$[2023Q1] & & & -0.01078 \\
\midrule
Other characteristics & Yes (9) & Yes (9) & Yes (9) \\
Observations & 960 & 953 & 1,913 \\
Unique banks & 960 & 953 & 994 \\
$R^2$ & 0.2920 & 0.0782 & --- \\
Standard errors & Classical & Classical & Clustered (bank) \\
\bottomrule
\end{tabular}
\begin{minipage}{0.95\textwidth}
\vspace{2mm}\footnotesize
\textit{Notes.} Standard errors in parentheses, $p$-values in brackets. The dependent variable is the quarterly deposit growth rate. All ten balance-sheet characteristics, the sample screens (assets $>$ \$1bn, deposits $\geq$ 50\% of liabilities, the FDIC charter screen) and the specification replicate the main analysis exactly, with the dates shifted: in the pre-crisis comparison period every predictor is measured at 2022Q1 and the outcome runs 2022Q1 to 2022Q2, so no post-period information enters. The FDIC failed-bank list records no failures in 2021 or 2022, so nothing is censored in the comparison period; the banking-stress column retains the baseline treatment of the three 2023 failures (censored at $-1.0$). The pooled column stacks both periods' bank-period observations, each having passed its own screens, and clusters standard errors on the bank because 919 banks appear in both periods. The interaction coefficient is the test of whether the relationship is specific to the stress episode.
\end{minipage}
\end{table}
